\begin{landscape}
\begingroup
\fontsize{10}{13}\selectfont
\fancyhead[L]{\small POSTHUMOUS PAPERS. PERTURBATIONS OF PALLAS}
\renewcommand{\headrulewidth}{0pt}
\par\medskip
\SourceLeaf{563}
\pdfbookmark[2]{Posthumous papers VI [9]}{pallas-VI-fragment9}
\begin{center}[9.]\end{center}
\makeatletter
\providecommand{\VInineHistoricalNote}[1]{\begingroup
\renewcommand{\thefootnote}{*)}
\long\def\@makefntext##1{\noindent [*)\ ##1}
\footnote{#1}\endgroup}
\makeatother
\par\noindent\hspace*{\parindent}[As already noted in Article 19 of the Exposition, Pallas provides a case in which the mean motions \(n\) and \(n'\) are in a rational ratio. The mean motions of Pallas and Jupiter are in the exact ratio \(18:7\), which is always exactly restored\VInineHistoricalNote{See \textsc{Gauss} to \textsc{Bessel}, 5 May 1812, p. 421 of this volume, and Göttingische gelehrte Anzeigen for 25 April 1812, vol. VI, p. 349.]}. Here the terms depending on the argument \(18\text{\mdseries\AstroSymbols ♃}-7\text{\mdseries\AstroSymbols ⚴}\) are to be given in secular form for practical calculation. The next article will treat the exact integration by which their periodic form is recovered.
\par\noindent\hspace*{\parindent}From the numerical values of \(T,V,W\) above, after interpolation as in Articles 2—7, we obtain the terms depending on that argument
\[
\begin{aligned}
\text{in }\frac{(1-\cos i)\mathrm{d}\AscendingNode}{n\mathrm{d}t}:&\quad
0\overset{\prime\prime}{,}000\cos(18M'-7M)\\
&\qquad+0\overset{\prime\prime}{,}001\sin(18M'-7M)\\[.7ex]
\text{in }\frac{\mathrm{d}\tilde\omega}{n\mathrm{d}t}-\frac{(1-\cos i)\mathrm{d}\AscendingNode}{n\mathrm{d}t}:&\quad
+0\overset{\prime\prime}{,}095\cos(18M'-7M)\\
&\qquad-0\overset{\prime\prime}{,}055\sin(18M'-7M)\\[.7ex]
\text{in }\frac{\mathrm{d}\log\mathrm{hyp}\,n}{n\mathrm{d}t}=-\tfrac32\frac{\mathrm{d}\log\mathrm{hyp}\,a}{n\mathrm{d}t}:&\quad
+0\overset{\prime\prime}{,}012\cos(18M'-7M)\\
&\qquad-0\overset{\prime\prime}{,}027\sin(18M'-7M)\\[.7ex]
\text{in }\frac{\mathrm{d}\varepsilon}{n\mathrm{d}t}-\frac{(1-\cos i)\mathrm{d}\AscendingNode}{n\mathrm{d}t}:&\quad
-0\overset{\prime\prime}{,}039\cos(18M'-7M)\\
&\qquad-0\overset{\prime\prime}{,}021\sin(18M'-7M).
\end{aligned}
\]
\par\noindent\hspace*{\parindent}Combining the two terms in \(\cos(18M'-7M)\) and \(\sin(18M'-7M)\) into one and introducing the mean longitudes through the values
\[
\begin{aligned}
M'&=\text{\mdseries\AstroSymbols ♃}-11^0 17'\quad5\overset{\prime\prime}{,}39\\
M&=\text{\mdseries\AstroSymbols ⚴}-121\quad8\quad54{,}50,
\end{aligned}
\]
gives the term
\[
\begin{aligned}
\text{in }\tilde\omega:&\quad-0\overset{\prime\prime}{,}1093\int\cos(18\text{\mdseries\AstroSymbols ♃}-7\text{\mdseries\AstroSymbols ⚴}+134^0 31'\,38'')\mathrm{d}\text{\mdseries\AstroSymbols ⚴}\\
\text{in }\log a:&\quad+0{,}415\int\cos(18\text{\mdseries\AstroSymbols ♃}-7\text{\mdseries\AstroSymbols ⚴}+170^0 56'\,57'')\mathrm{d}\text{\mdseries\AstroSymbols ⚴}\\
&\qquad\text{(in units of the 7th decimal place)}\\
\text{in }\int n\mathrm{d}t:&\quad-0\overset{\prime\prime}{,}02955\iint\cos(18\text{\mdseries\AstroSymbols ♃}-7\text{\mdseries\AstroSymbols ⚴}+170^0 56'\,57'')\mathrm{d}\text{\mdseries\AstroSymbols ⚴}^{2}\\
\text{in }\varepsilon:&\quad-0\overset{\prime\prime}{,}0438\int\cos(18\text{\mdseries\AstroSymbols ♃}-7\text{\mdseries\AstroSymbols ⚴}+257^0 45'\,45'')\mathrm{d}\text{\mdseries\AstroSymbols ⚴};
\end{aligned}
\]
Integration, with \(18\text{\mdseries\AstroSymbols ♃}-7\text{\mdseries\AstroSymbols ⚴}\) set equal to the constant value \(125^0 36'\,49''\), gives:
\begin{center}
\begin{tabular}{@{}cccll@{}}
Secular change&of&\(\tilde\omega\)&\(:+0\overset{\prime\prime}{,}0255\)&(annually)\\
»&»&\(\log a\)&\(:+0{,}25\)&»\\
&&\multicolumn{3}{l}{(in units of the 7th decimal place)}\\
»&»&\(\varepsilon\)&\(:-0\overset{\prime\prime}{,}0548\)&»\\
\multicolumn{3}{r}{Coefficient of \(tt\) in \(\int n\mathrm{d}t\)}&\(:-0\overset{\prime\prime}{,}0123\)&»
\end{tabular}
\end{center}
\par\noindent\hspace*{\parindent}For the inclination and the longitude of the node, the corresponding terms are very small.]
\clearpage
\endgroup
\end{landscape}
